\documentclass[11pt]{article}
\usepackage{amsmath,amssymb,amsthm}
\usepackage{mathtools}
\usepackage[margin=1in]{geometry}
\usepackage{booktabs}
\usepackage{hyperref}
\usepackage{xcolor}
\hypersetup{colorlinks=true, linkcolor=blue!50!black, urlcolor=blue!50!black}

\newcommand{\Var}{\mathcal{V}}
\newcommand{\code}[1]{\texttt{#1}}
\DeclareMathOperator{\Sol}{Sol}
\DeclareMathOperator{\GCRD}{GCRD}
\DeclareMathOperator{\ord}{ord}
\newcommand{\Vall}{V_{\forall}}
\newcommand{\Vex}{V_{\exists\backslash\Psi}}
\newcommand{\vlin}{v_{\mathrm{lin}}}

\newtheorem{theorem}{Theorem}
\newtheorem{lemma}{Lemma}
\newtheorem{proposition}{Proposition}
\newtheorem{corollary}{Corollary}
\newtheorem{definition}{Definition}
\newtheorem{remark}{Remark}

\title{When the membership and homogeneous-existence loci coincide,\\
       and a proof that they do for ansatz~5 on hydrogen}
\author{}
\date{\today}

\begin{document}
\maketitle

\noindent
A companion note to \emph{A New Method of Solving PDEs}, and a sequel to
\emph{Membership versus variety, and the partial-solution strata} and
\emph{Recovering the partial-solution strata by membership alone: a GCRD closure
construction}. The core algorithm computes the \textbf{membership locus}
$\Vall(A)=\{c:P\in[A(c)]\}$; the existential computation on $A\cup\{P\}$ computes
the \textbf{homogeneous existence locus} $\Vex(A)=\{c:1\notin[A(c),P]:\Psi^\infty\}$.
This note (i) records the exact condition under which the two coincide, (ii)
proves coincidence unconditionally for one-dimensional ans\"atze, and (iii)
proves by an explicit greatest-common-right-divisor (GCRD) computation that
\emph{ansatz~5 on the hydrogen Schr\"odinger equation has no genuine partial
strata}, so that the two loci coincide there off the bad locus, and the whole
gap between them is confined to the ODE-degeneration loci that the differential
Thomas decomposition already resolves cell by cell.

\section{The exact gap, and its GCRD reading}

The relationship between the two loci was established in the two companion
notes. Writing the existential projection as membership plus a remainder,
\begin{equation}
\label{eq:gap}
\Vex(A) \;=\; \Vall(A)\ \sqcup\ \{\text{partial-solution strata}\},
\end{equation}
where a \emph{partial-solution stratum} (companion note~I, Definition~1) is a
$c$ at which a proper nontrivial subfamily of $\Sol(A(c))$ satisfies $P$ while
the family as a whole does not. Hence
\begin{equation}
\label{eq:coincide}
\boxed{\ \Vall(A)=\Vex(A)\ \iff\ A\ \text{has no partial-solution strata.}\ }
\end{equation}

For an ansatz whose top element is a \emph{linear} ODE $L$ of order $n$ in the
inner variable $v$, companion note~II sharpens this into an order condition.
Ritt-reducing $P$ modulo $A$ and splitting the remainder over base power
products yields operators $R_1,\dots,R_s\in K(v)[\partial]$ of order $\le n-1$;
with
\[
G(c^*)\;=\;\GCRD\big(L(c^*),R_1(c^*),\dots,R_s(c^*)\big),\qquad k=\ord G(c^*),
\]
the trichotomy (note~II, Theorem~1) reads: $k=n\iff c^*\in\Vall$ (the whole
family solves $P$); $k=0\iff$ no common solution; and $0<k<n\iff c^*$ is a
partial stratum. Therefore
\begin{equation}
\label{eq:gcrd-coincide}
\Vall(A)=\Vex(A)\quad\text{off the bad locus}\quad\iff\quad
\ord G(c^*)\in\{0,n\}\ \text{ for all such }c^* .
\end{equation}
Coincidence is exactly the statement that the GCRD never has intermediate
order.

\begin{remark}[hypotheses and their necessity]
\label{rem:hyp}
Everything past \eqref{eq:coincide} rests on two standing hypotheses, and both
are load-bearing.

\emph{(H1) $P$ is linear and homogeneous in $\Psi$.} Linearity is what makes
the reduced remainder a \emph{linear} operator: $r=\sum_\beta m_\beta
R_\beta(\Psi)$ with $R_\beta\in K(v)[\partial]$, the object whose GCRD with $L$
the whole analysis takes. A $\Psi^2$- or $|\Psi|^2\Psi$-term would make $r$
nonlinear in the jets, and there would be no operator to divide --- the Ore
ring, the kernel-intersection identity, the Wronskian argument and the
subresultant strata equations all require $r$ to be an operator. Homogeneity is
what makes the $P$-solving solutions a linear \emph{subspace} of
$\Sol(A(c^*))$, on which Lemma~\ref{lem:dim1} turns; without it a
one-dimensional family could solve $P$ at isolated scalings and the dimension
argument would fail. Linearity here means linear in $\Psi$ and its
derivatives; arbitrary coefficients in the independent variables are allowed,
so the Coulomb $-1/r$ is fine and the Schr\"odinger equation satisfies~(H1). A
nonlinear Schr\"odinger / Gross--Pitaevskii equation does not, and lies outside
the theory of this note entirely.

\emph{(H2) the ansatz top element is a linear ODE.} This is what places $L$ in
the right-Euclidean ring $K(v)[\partial]$, where GCRDs exist and are
computable; the trichotomy is a statement about $\ord\GCRD$ of linear
operators. Nonlinear structure \emph{below} the top is admissible provided it
is absorbed into the differential base field: the algebraic-extension elements
of ans\"atze 13 and 16, the roots, and the coefficient-ring extensions 14, 15
leave the top element linear (over $K(v)(g)$, say), and GCRD theory is valid
over any differential field. Companion note~II further relaxes~(H2) for a
\emph{first-order} nonlinear top, which has an algebraic-extension closure of
bounded degree; a \emph{higher-order} nonlinear top is open. The one library
ansatz that genuinely violates~(H2) is the product family~10
($\Psi=F(\xi)\,G(\eta)$): eliminating $F,G$ from the nonlinear relation
$\Psi-F\!\cdot\!G$ gives a nonlinear top for $\Psi$, outside $K(v)[\partial]$.

Both hydrogen and helium are linear and homogeneous in $\Psi$, so~(H1) holds
throughout; every library ansatz except~10 satisfies~(H2). The
ansatz-5 result of \S\ref{sec:ansatz5} sits squarely in the doubly-linear
regime.
\end{remark}

\section{One-dimensional ans\"atze coincide unconditionally}

The cleanest sufficient condition is dimensional and needs no GCRD machinery.

\begin{lemma}[dimension-one coincidence]
\label{lem:dim1}
Let $P$ be linear and homogeneous in $\Psi$. If, for every $c^*$ off the bad
locus, $\dim_{\mathbb C}\Sol(A(c^*))\le 1$, then $\Vall(A)=\Vex(A)$ off the bad
locus.
\end{lemma}

\begin{proof}
For $P$ linear and homogeneous, $\{\Psi\in\Sol(A(c^*)):P(\Psi)=0\}$ is a
\emph{linear} subspace of $\Sol(A(c^*))$. A space of dimension $\le 1$ has only
$\{0\}$ and itself as subspaces, so either every ansatz solution solves $P$
(then $c^*\in\Vall$) or only the zero solution does (then $c^*\notin\Vex$). In
both cases $c^*$ lies in $\Vall$ iff it lies in $\Vex$. In the GCRD reading,
$n\le 1$ forces $k\in\{0,n\}$, so \eqref{eq:gcrd-coincide} holds.
\end{proof}

The Schr\"odinger operator is linear and homogeneous in $\Psi$, so
Lemma~\ref{lem:dim1} applies to every ansatz of the library whose solution
space is one-dimensional --- the first-order-ODE and pure-exponential families.
The two-dimensional (second-order ODE) families are not covered, and, as
companion note~I shows with $A:\Psi''-c\Psi$, $P:\Psi''-3\Psi'+2\Psi$, they can
carry genuine partial strata even when the PDE order matches the ansatz order.
Whether a \emph{particular} second-order ansatz has strata against a
\emph{particular} PDE is the residual question; \S\ref{sec:ansatz5} settles it
for ansatz~5 on hydrogen.

\begin{center}
\begin{tabular}{@{}lll@{}}
\toprule
$\dim\Sol(A(c))$ & library ans\"atze & status \\
\midrule
$1$ (1st-order ODE / pure exp) & 8, 9, 1, 1.1, 20, 20.1 & coincide (Lemma~\ref{lem:dim1}) \\
$1$? (log additive mode) & 2, 17, 17.1, 18, 19 & dimension check pending \\
$2$ (2nd-order linear ODE) & 5, 5.1, 5.2, 5.3, 3, 6, 7, 13, 14, 15, 16 & GCRD analysis (\S\ref{sec:ansatz5}) \\
higher / nested & 12 & GCRD tower recursion \\
nonlinear (product $F\!\cdot\!G$) & 10 & outside the linear theory \\
\bottomrule
\end{tabular}
\end{center}

\section{Ansatz~5 on hydrogen has no genuine partial strata}
\label{sec:ansatz5}

We carry out the GCRD analysis of \eqref{eq:gcrd-coincide} for the paper's
figure ``ansatz~5'' against the hydrogen Schr\"odinger equation, and find that
the only partial strata sit on the bad locus. The computation is elementary
linear algebra over $\mathbb{Q}(c)$; the verification scripts are
\code{ansatz5\_gcrd.sage} and \code{ansatz5\_gcrd2.sage}.

\subsection{Setup and the reduced remainder}

Ansatz~5 is $\Psi=\zeta(v)$ with inner variable
$v=v_1x+v_2y+v_3z+v_4r$, $r^2=x^2+y^2+z^2$, and second-order top element
\begin{equation}
\label{eq:L}
L\;=\;A_2\,\partial^2+L_1\,\partial+L_0,\qquad
A_2=a_0+a_1v,\ L_1=b_0+b_1v,\ L_0=c_0+c_1v,\ \partial=\tfrac{d}{dv}.
\end{equation}
Write $\vlin=(v_1,v_2,v_3)$ and $K=|\vlin|^2-v_4^2$. Since
$v_{\mathrm{lin}}\!\cdot x=v-v_4r$, the gradient and Laplacian of the inner
variable are
\[
|\nabla v|^2=K+\frac{2v_4v}{r},\qquad \nabla^2v=\frac{2v_4}{r},
\]
so, for $\Psi=\zeta(v)$,
$\nabla^2\Psi=\zeta''\big(K+\tfrac{2v_4v}{r}\big)+\zeta'\,\tfrac{2v_4}{r}$.
Clearing the Schr\"odinger equation $-\tfrac12\nabla^2\Psi-\tfrac1r\Psi-E\Psi=0$
by $-2r$ gives
\begin{equation}
\label{eq:prereduce}
\zeta''\,(Kr+2v_4v)+2v_4\,\zeta'+2\zeta+2Er\,\zeta \;=\;0 .
\end{equation}
Reducing $\zeta''$ through $L$ --- multiply \eqref{eq:prereduce} by $A_2$ and
substitute $A_2\zeta''=-L_1\zeta'-L_0\zeta$ --- and collecting on the
\emph{independent} radical $r$ (valid where $|\vlin|\neq0$, so that $v$ and $r$
vary independently; Remark~\ref{rem:generic}) leaves a remainder
\begin{equation}
\label{eq:rem}
r\,\big(R_1\zeta\big)+\big(R_0\zeta\big)\;=\;0,
\qquad h=A_2=a_0+a_1v,
\end{equation}
where the two first-order remainder operators are
\begin{align}
\label{eq:R1}
R_1&=-K L_1\,\partial+\big(2E A_2-K L_0\big),\\
\label{eq:R0}
R_0&=2v_4\,(A_2-vL_1)\,\partial+2\,(A_2-v_4vL_0).
\end{align}
Because $v$ and $r$ are independent, \eqref{eq:rem} forces $R_1\zeta=0$ and
$R_0\zeta=0$ separately. A partial stratum off the bad locus ($h\ne0$) is thus a
parameter point at which $L$, $R_1$, $R_0$ share a nonzero solution, i.e.\
$\ord\GCRD(L,R_1,R_0)=1$.

\begin{remark}[base-splitting genericity]
\label{rem:generic}
The split on $r$ requires $v$ and $r$ to be independent coordinates, i.e.\
$|\vlin|\neq0$; it also uses $A_2\ne0$ (the multiplier $h$) and, to keep $R_1$
of order one, $K\neq0$ and $v_4\neq0$. The excluded loci $\{|\vlin|=0\}$
(radial inner variable, $v=v_4r$), $\{K=0\}$, $\{v_4=0\}$ and $\{h=0\}$ are
precisely the ODE-degeneration strata handled cell by cell by the differential
Thomas decomposition; the radial locus $\{|\vlin|=0\}$ is where the hydrogen
bound states live.
\end{remark}

\subsection{The GCRD computation}

With $\rho=-B_1/A_1$ the exponent of the $R_1$-mode ($A_1=-KL_1$,
$B_1=2EA_2-KL_0$), a common solution of $L,R_1,R_0$ exists iff
\begin{itemize}
\item[(P)] $R_0$ selects the same mode: $A_1B_0-A_0B_1=0$ identically in $v$; and
\item[(R)] that mode solves $L$: $A_2(\rho'+\rho^2)+L_1\rho+L_0=0$ (the Riccati
of $L$), cleared by $A_1^2$ to
$A_2\big(B_1^2-B_1'A_1+B_1A_1'\big)-L_1B_1A_1+L_0A_1^2=0$.
\end{itemize}
Each is an identity in $v$; equating $v$-coefficients gives polynomial equations
in $(a_0,a_1,b_0,b_1,c_0,c_1,K,v_4,E)$. Let $J$ be their ideal.

\begin{theorem}[no genuine partial strata for ansatz~5 / hydrogen]
\label{thm:ansatz5}
Off the degenerate loci $\{K=0\}\cup\{v_4=0\}\cup\{|\vlin|=0\}$, the ansatz~5
system for the hydrogen Schr\"odinger equation has no partial-solution stratum
with $h\not\equiv 0$. Consequently $\Vall=\Vex$ off the bad locus, and the
entire gap $\Vex\setminus\Vall$ lies on the ODE-degeneration loci.
\end{theorem}

\begin{proof}
Saturating $J$ by $K$ and $v_4$ (genericity) and then by $(a_0,a_1)$
(imposing $h\not\equiv 0$) and computing a Gr\"obner basis, one finds the basis
contains $b_0^2,\ b_0b_1,\ b_1^2$; hence every solution with $h\not\equiv 0$
lies on $\{b_0=b_1=0\}$, i.e.\ $L_1\equiv0$. Saturating additionally by
$(b_0,b_1)$ --- the branch $L_1\neq0$ --- yields the unit ideal: there is no
partial stratum with a genuine first-derivative term.

On the residual branch $L_1\equiv0$ the derivation above degenerates, since
$A_1=-KL_1=0$ and $R_1$ drops to order zero; it must be treated directly.
There $R_1=B_1$, so $R_1\zeta=0$ forces $B_1=0$, i.e.\ $2EA_2=KL_0$, whence
$L_0=(2E/K)A_2$ and, dividing $L$ by $A_2\neq0$, the ODE is
$\zeta''+(2E/K)\zeta=0$. The remaining operator is
$R_0=2A_2\big(v_4\partial+1-(2Ev_4/K)v\big)$, giving
$\rho_0=\zeta'/\zeta=-1/v_4+(2E/K)v$; requiring $\rho_0$ to solve the ODE's
Riccati $\rho_0'+\rho_0^2+2E/K=0$ and clearing by $(Kv_4)^2$ gives
\[
4v_4^2E^2\,v^2-4Kv_4E\,v+4Kv_4^2E+K^2\;=\;0\quad(\text{in }v),
\]
whose $v^2$-coefficient forces $E=0$ and whose constant term then forces
$K^2=0$, i.e.\ $K=0$, excluded. So this branch is empty as well. Both branches
being empty, no $h\not\equiv0$ partial stratum exists; by \eqref{eq:coincide}
and \eqref{eq:gcrd-coincide}, $\Vall=\Vex$ off the bad locus, and
\eqref{eq:gap} leaves the gap on the degeneration loci.
\end{proof}

\subsection{Reading}

For ansatz~5 on hydrogen the second-order solution space never lets the
Schr\"odinger equation peel off a single mode away from the degeneration loci:
the potential's $-1/r$ term and the $K,v_4$ geometry conspire so that (P) and
(R) are jointly unsatisfiable with $h\neq0$. Consequently the membership
computation, run \emph{per Thomas cell}, already recovers the full homogeneous
existence locus: the cells on the bad locus $\{h=0\}$ and the radial locus
$\{|\vlin|=0\}$ --- where the bound-state energies $E=-\tfrac12,-\tfrac18$ sit
--- are the only places $\Vall$ and $\Vex$ can differ, and those are exactly the
degenerate cells the decomposition isolates. No separate GCRD $\forall$-cover
(companion note~II) is needed for ansatz~5; it remains the general safety net
for ans\"atze whose GCRD does reach intermediate order.

\begin{remark}[scope]
Theorem~\ref{thm:ansatz5} is specific to the \emph{hydrogen} Hamiltonian and to
ansatz~5's linear-coefficient second-order ODE. It does not follow from
order-matching (companion note~I gives a both-second-order counterexample with
genuine strata) and does not transfer automatically: the helium Hamiltonian,
and the quadratic-coefficient variants 5.1--5.3 and the coefficient-ring
extensions 13--16, require re-running the same computation. The scripts
\code{ansatz5\_gcrd.sage}/\code{ansatz5\_gcrd2.sage} are parameterized to make
that a mechanical substitution.
\end{remark}

\section{Summary}

\begin{itemize}
\item $\Vall=\Vex\iff$ no partial strata $\iff$ (linear-ODE ans\"atze)
      $\ord\GCRD(L,\{R_\beta\})\in\{0,n\}$ everywhere (\eqref{eq:coincide},
      \eqref{eq:gcrd-coincide}).
\item One-dimensional ans\"atze coincide unconditionally
      (Lemma~\ref{lem:dim1}); this covers the first-order and pure-exponential
      library families.
\item Ansatz~5 on hydrogen has no genuine ($h\neq0$) partial strata
      (Theorem~\ref{thm:ansatz5}); its $\Vall$-vs-$\Vex$ gap is entirely on the
      ODE-degeneration loci resolved by the Thomas decomposition, so per-cell
      membership recovers $\Vex$ with no existential computation.
\item Open: the log ans\"atze dimension check (2, 17, 17.1, 18, 19); the
      quadratic-coefficient and coefficient-ring second-order families; the
      nested (12) and product (10) ans\"atze; and the helium Hamiltonian.
\end{itemize}

\end{document}
